\setcounter{section}{13}






  \zwischenueberschrift{Übungsaufgaben}

\inputaufgabe
{}
{





Es sei $R$ ein
\definitionsverweis {kommutativer Ring}{}{}
und

\mavergleichskette
{\vergleichskette
{ S
}
{ \subseteq }{ R
}
{  }{
}
{    }{
}
{   }{
}
}
{}{}{}
ein
\definitionsverweis {multiplikatives System}{}{.}
Man definiert die
\definitionswortenp{Nenneraufnahme}{}

\mathdisp {R_S} {  }

schrittweise wie folgt. Es sei zunächst $M$ die Menge der formalen Brüche mit Nenner in $S$, also

\mavergleichskettedisp
{\vergleichskette
{M
}
{ =} { \left\{  \frac{r}{s}   \mid   r \in R , \,  s \in S       \right\}
}
{ } {
}
{ } {
}
{ } {
}
}
{}{}{.}
Zeige, dass durch

\mathdisp {\frac{r}{s} \sim \frac{r'}{s'}  \text{ genau dann, wenn es ein } t \in S   \text{ mit }    trs' =tr's  \text{ gibt}} { , }

eine
\definitionsverweis {Äquivalenzrelation}{}{}
auf $M$ definiert ist. Wir bezeichnen mit $R_S$ die Menge der Äquivalenzklassen. Definiere auf $R_S$ eine Ringstruktur und definiere einen
\definitionsverweis {Ringhomomorphismus}{}{}

\mathl{R \rightarrow R_S}{.}





}
{} {}

\inputaufgabe
{}
{





Es sei $R$ ein
\definitionsverweis {Integritätsbereich}{}{}
und sei

\mavergleichskette
{\vergleichskette
{S
}
{ \subseteq }{ R
}
{  }{
}
{    }{
}
{   }{
}
}
{}{}{}
ein
\definitionsverweis {multiplikatives System}{}{,}

\mavergleichskette
{\vergleichskette
{ 0
}
{ \notin }{S
}
{  }{
}
{    }{
}
{   }{
}
}
{}{}{.}
\aufzaehlungzwei {Zeige, dass die
\definitionsverweis {Nenneraufnahme}{}{}
zu $S$, also $R_S$ mit

\mavergleichskettedisp
{\vergleichskette
{R_S
}
{ \defeq} { \left\{  \frac{f}{g}  \mid   f \in R  , \,  g \in S       \right\}
}
{ \subseteq} { Q(R)
}
{ } {
}
{ } {
}
}
{}{}{}
ein Unterring von
\mathl{Q(R)}{} ist.
} {Zeige, dass nicht jeder Unterring von $Q(R)$ eine Nenneraufnahme ist.}





}
{} {}

\inputaufgabegibtloesung
{}
{





Zeige, dass der Körper der
\definitionsverweis {rationalen Zahlen}{}{}
$\Q$
\definitionsverweis {überabzählbar}{}{}
viele
\definitionsverweis {Unterringe}{}{}
besitzt.





}
{} {}

\inputaufgabe
{}
{





Es sei $R$ ein
\definitionsverweis {kommutativer Ring}{}{}
und sei

\mavergleichskette
{\vergleichskette
{ f
}
{ \in }{ R
}
{  }{
}
{    }{
}
{   }{
}
}
{}{}{}
mit zugehöriger
\definitionsverweis {Nenneraufnahme}{}{}
$R_f$. Beweise die
$R$-\definitionsverweis {Algebraisomorphie}{}{}

\mavergleichskettedisp
{\vergleichskette
{R_f
}
{ \cong} { R[T]/(Tf- 1)
}
{ } {
}
{ } {
}
{ } {
}
}
{}{}{.}





}
{} {}

\inputaufgabe
{}
{





Es sei $R$ ein
\definitionsverweis {kommutativer Ring}{}{,}

\mavergleichskette
{\vergleichskette
{ f
}
{ \in }{ R
}
{  }{
}
{    }{
}
{   }{
}
}
{}{}{}
ein Element und $R_f$ die zugehörige
\definitionsverweis {Nenneraufnahme}{}{.}
Zeige, dass $f$ genau dann
\definitionsverweis {nilpotent}{}{}
ist, wenn $R_f$ der Nullring ist.





}
{} {}

In den folgenden Aufgaben dürfen Sie, wenn Sie wollen, bei Nenneraufnahmen annehmen, dass Integritätsbereiche vorliegen.

\inputaufgabegibtloesung
{}
{





Es seien $R$ und $A$
\definitionsverweis {kommutative Ringe}{}{}
und sei

\mavergleichskette
{\vergleichskette
{ S
}
{ \subseteq }{ R
}
{  }{
}
{    }{
}
{   }{
}
}
{}{}{}
ein
\definitionsverweis {multiplikatives System}{}{.}
Es sei
\maabbdisp {\varphi} { R } { A
} {}
ein
\definitionsverweis {Ringhomomorphismus}{}{}
derart, dass $\varphi(s)$ eine
\definitionsverweis {Einheit}{}{}
in $A$ ist für alle

\mavergleichskette
{\vergleichskette
{ s
}
{ \in }{ S
}
{  }{
}
{    }{
}
{   }{
}
}
{}{}{.}
Zeige: Dann gibt es einen eindeutig bestimmten Ringhomomorphismus
\maabbdisp {\tilde{\varphi}} { R_S } { A
} {,}
der $\varphi$ fortsetzt.





}
{} {}

\inputaufgabe
{}
{





Es sei $R$ ein
\definitionsverweis {kommutativer Ring}{}{}
und

\mavergleichskette
{\vergleichskette
{ S
}
{ \subseteq }{ R
}
{  }{
}
{    }{
}
{   }{
}
}
{}{}{}
ein
\definitionsverweis {multiplikatives System}{}{.}
Zeige, dass die
\definitionsverweis {Primideale}{}{}
in $R_S$ genau denjenigen Primidealen in $R$  entsprechen, die mit $S$  einen leeren Durchschnitt haben.





}
{} {}

\inputaufgabegibtloesung
{}
{





Es sei $K$ ein
\definitionsverweis {Körper}{}{,}
$R=K[X,Y]$ der
\definitionsverweis {Polynomring}{}{}
in zwei Variablen, $S \subseteq R$ ein
\definitionsverweis {multiplikatives System}{}{}
und $F \in R$ ein Polynom. Zeige, dass es eine eindeutige
$R$-\definitionsverweis {Algebraisomorphie}{}{}

\mavergleichskettedisp
{\vergleichskette
{ (R/(F))_S
}
{ \cong} { (R_S)/(F)
}
{ } {
}
{ } {
}
{ } {
}
}
{}{}{}
gibt.





}
{} {}

\inputaufgabegibtloesung
{}
{





Es sei $R$ ein
\definitionsverweis {kommutativer Ring}{}{,}

\mavergleichskette
{\vergleichskette
{ {\mathfrak a}
}
{ \subseteq }{ R
}
{  }{
}
{    }{
}
{   }{
}
}
{}{}{}
ein
\definitionsverweis {Ideal}{}{}
und

\mavergleichskette
{\vergleichskette
{S
}
{ \subseteq }{ R
}
{  }{
}
{    }{
}
{   }{
}
}
{}{}{}
ein
\definitionsverweis {multiplikatives System}{}{.}
Zeige, dass es eine natürliche
\definitionsverweis {Ringisomorphie}{}{}

\mavergleichskettedisp
{\vergleichskette
{ \left(  R/ {\mathfrak a}  \right)_S
}
{ \cong} { R_S/ {\mathfrak a} R_S
}
{ } {
}
{ } {
}
{ } {
}
}
{}{}{}
gibt, wobei links die
\definitionsverweis {Nenneraufnahme}{}{}
am Bild des multiplikativen Systems in $R/ {\mathfrak a}$ bezeichnet.





}
{} {}

\inputaufgabe
{}
{





Es sei $K$ ein
\definitionsverweis {algebraisch abgeschlossener Körper}{}{}
und seien $R$ und $S$ kommutative
$K$-\definitionsverweis {Algebren von endlichem Typ}{}{.}
Es sei

\mavergleichskette
{\vergleichskette
{ f
}
{ \in }{ R
}
{  }{
}
{    }{
}
{   }{
}
}
{}{}{}
und
\maabb {\varphi} { R} { S
} {}
sei ein
$K$-\definitionsverweis {Algebrahomomorphismus}{}{.}
Zeige, dass die
\definitionsverweis {Spektrumsabbildung}{}{}
$\varphi^*$ genau dann durch $D(f)$ faktorisiert, wenn $\varphi(f)$ eine
\definitionsverweis {Einheit}{}{}
in $S$ ist.





}
{} {}

\inputaufgabegibtloesung
{}
{





Es sei $K$ ein
\definitionsverweis {algebraisch abgeschlossener Körper}{}{}
und $R$ eine
\definitionsverweis {integre}{}{}
\definitionsverweis {endlich erzeugte}{}{}
$K$-\definitionsverweis {Algebra}{}{.}
Es seien

\mavergleichskette
{\vergleichskette
{ f,g
}
{ \in }{ R
}
{  }{
}
{    }{
}
{   }{
}
}
{}{}{.}
Zeige, dass die folgenden Aussagen äquivalent sind:
\aufzaehlungzwei {
\mavergleichskette
{\vergleichskette
{ D(f)
}
{ \subseteq }{ D(g)
}
{  }{
}
{    }{
}
{   }{
}
}
{}{}{.}
} {Es gibt einen
$R$-\definitionsverweis {Algebrahomomorphismus}{}{}
\maabb {} {R_g } { R_f
} {.}
}
Zeige ferner, dass diese Äquivalenz für

\mavergleichskette
{\vergleichskette
{ K
}
{ = }{ \R
}
{  }{
}
{    }{
}
{   }{
}
}
{}{}{}
nicht gilt.





}
{} {}

Die folgende Aufgabe verwendet den Begriff des saturierten multiplikativen Systems.





Ein \definitionsverweis {multiplikatives System}{}{}
$S$ in einem
\definitionsverweis {kommutativen Ring}{}{}
$R$ heißt \definitionswort {saturiert}{,} wenn Folgendes gilt: Ist

\mavergleichskette
{\vergleichskette
{ g
}
{ \in }{ R
}
{  }{
}
{    }{
}
{   }{
}
}
{}{}{}
und gibt es ein

\mavergleichskette
{\vergleichskette
{ f
}
{ \in }{ S
}
{  }{
}
{    }{
}
{   }{
}
}
{}{}{,}
das von $g$ geteilt wird, so ist auch

\mavergleichskette
{\vergleichskette
{ g
}
{ \in }{ S
}
{  }{
}
{    }{
}
{   }{
}
}
{}{}{.}







\inputaufgabe
{}
{





Es seien $A,B$
\definitionsverweis {kommutative Ringe}{}{}
und sei
\mathl{\varphi:A \rightarrow B}{} ein
\definitionsverweis {Ringhomomorphismus}{}{.}
Zeige, dass das Urbild
\mathl{\varphi^{-1}(B^\times)}{} der Einheitengruppe ein
\definitionsverweis {saturiertes multiplikatives System}{}{}
in $A$ ist.





}
{} {}

\inputaufgabe
{}
{





Es sei $R$ ein
\definitionsverweis {kommutativer Ring}{}{.}
Zeige, dass die Menge der
\definitionsverweis {Nichtnullteiler}{}{}
in $R$ ein
\definitionsverweis {saturiertes}{}{}
\definitionsverweis {multiplikatives System}{}{}
bildet.





}
{} {}

\inputaufgabegibtloesung
{}
{





Man gebe ein Beispiel einer integren, endlich erzeugten ${\mathbb C}$-Algebra $R$ und eines multiplikativen Systems
\mathbed {S \subseteq R} {}
 {0 \not\in S} {}
 {} {} {} {,} an derart, dass die Nenneraufnahme $R_S$ kein Körper ist, aber jedes maximale Ideal aus $R$ zum Einheitsideal in $R_S$ wird.





}
{} {}

\inputaufgabegibtloesung
{}
{





Zeige, dass ein \definitionsverweis {Integritätsbereich}{}{}
ein
\definitionsverweis {zusammenhängender Ring}{}{}
ist.





}
{} {}

\inputaufgabe
{}
{





Es sei $R$ ein kommutativer Ring und sei

\mavergleichskette
{\vergleichskette
{ f
}
{ \in }{ R
}
{  }{
}
{    }{
}
{   }{
}
}
{}{}{.}
Es sei $f$ sowohl
\definitionsverweis {nilpotent}{}{}
als auch
\definitionsverweis {idempotent}{}{.}
Zeige, dass

\mavergleichskette
{\vergleichskette
{ f
}
{ = }{ 0
}
{  }{
}
{    }{
}
{   }{
}
}
{}{}{}
ist.





}
{} {}

\inputaufgabegibtloesung
{}
{





Man gebe zu jedem

\mavergleichskette
{\vergleichskette
{ n
}
{ \geq }{ 2
}
{  }{
}
{    }{
}
{   }{
}
}
{}{}{}
einen
\definitionsverweis {kommutativen Ring}{}{}
$R$ und ein Element

\mathbed {x \in R} {}
 {x \neq 0} {}
 {} {} {} {,}
an, für das
\mathkor {} {nx=0} {und} {x^n =0} {}
gilt.





}
{} {}

\inputaufgabe
{}
{





Es sei $R$ ein
\definitionsverweis {kommutativer Ring}{}{}
und sei

\mavergleichskette
{\vergleichskette
{e
}
{ \in }{ R
}
{  }{
}
{    }{
}
{   }{
}
}
{}{}{}
ein
\definitionsverweis {idempotentes Element}{}{.}
Zeige, dass es eine natürliche
\definitionsverweis {Ringisomorphie}{}{}

\mavergleichskettedisp
{\vergleichskette
{ R_e
}
{ \cong} { R/(1-e)
}
{ } {
}
{ } {
}
{ } {
}
}
{}{}{}
gibt.





}
{\zusatzklammer {Dies zeigt erneut, dass $D(e)$ offen und abgeschlossen ist} {} {.}} {}

\inputaufgabe
{}
{





Es seien $R$ und $S$ kommutative Ringe und sei
\mathl{R \times S}{} der \definitionsverweis {Produktring}{}{}

\mathl{R \times S}{.} Zeige, dass die Teilmenge
\mathl{R \times 0}{} ein \definitionsverweis {Hauptideal}{}{}
ist.





}
{} {}

\inputaufgabegibtloesung
{}
{





Es sei

\mavergleichskette
{\vergleichskette
{ p
}
{ \in }{ \Z
}
{  }{
}
{    }{
}
{   }{
}
}
{}{}{}
eine Primzahl und

\mavergleichskette
{\vergleichskette
{ n
}
{ \in }{ \N
}
{  }{
}
{    }{
}
{   }{
}
}
{}{}{.}
Zeige, dass der
\definitionsverweis {Restklassenring}{}{}
$\Z/(p^n )$ nur die beiden trivialen
\definitionsverweis {idempotenten Elemente}{}{}
\mathkor {} {0} {und} {1} {}
besitzt.





}
{} {}

\inputaufgabegibtloesung
{}
{





Schreibe den
\definitionsverweis {Restklassenring}{}{}

\mathl{\Q[X]/ \left(  X^4-1  \right)}{} als ein Produkt von
\definitionsverweis {Körpern}{}{,}
wobei lediglich die Körper $\Q$ und $\Q[ { \mathrm i} ]$ vorkommen. Schreibe die Restklasse von $X^3+X$ als ein Tupel in dieser Produktzerlegung.





}
{} {}

\inputaufgabe
{}
{





Es sei $K$ ein
\definitionsverweis {Körper}{}{} und sei
\mathl{K[X]}{} der \definitionsverweis {Polynomring}{}{} über $K$. Es seien

\mavergleichskette
{\vergleichskette
{ a_1 , \ldots , a_n
}
{ \in }{ K
}
{  }{
}
{    }{
}
{   }{
}
}
{}{}{}
verschiedene Elemente und

\mavergleichskettedisp
{\vergleichskette
{ F
}
{ =} { (X-a_1){ \cdots }(X-a_n)
}
{ } {
}
{ } {
}
{ } {
}
}
{}{}{}
das Produkt der zugehörigen
\definitionsverweis {linearen Polynome}{}{.}
Zeige, dass der
\definitionsverweis {Restklassenring}{}{}
$K[X]/(F)$
\definitionsverweis {isomorph}{}{}
zum
\definitionsverweis {Produktring}{}{}
$K^n$ ist.





}
{} {}

\inputaufgabe
{}
{





Es sei $K$ ein
\definitionsverweis {algebraisch abgeschlossener Körper}{}{}
und
\mathl{K[X]}{} der
\definitionsverweis {Polynomring}{}{}
über $K$. Zeige, dass der
\definitionsverweis {Restklassenring}{}{}
zu einem Polynom

\mavergleichskette
{\vergleichskette
{F
}
{ \neq }{ 0
}
{  }{
}
{    }{
}
{   }{
}
}
{}{}{}
die Struktur

\mavergleichskettedisp
{\vergleichskette
{K[X]/(F)
}
{ \cong} { K[T]/ \left(  T^{n_1 }  \right) \times \cdots \times  K[T]/ \left(  T^{n_r}  \right)
}
{ } {
}
{ } {
}
{ } {
}
}
{}{}{}
besitzt. Zeige, dass dabei

\mavergleichskettedisp
{\vergleichskette
{ \operatorname{grad} \, ( F )
}
{ =} { n_1 + \cdots + n_r
}
{ } {
}
{ } {
}
{ } {
}
}
{}{}{}
ist.





}
{} {}

\inputaufgabegibtloesung
{}
{





Es sei $K$ ein
\definitionsverweis {Körper}{}{}
und seien

\mavergleichskette
{\vergleichskette
{ A
}
{ = }{ K[X_1  , \ldots , X_m ]/ {\mathfrak a}
}
{  }{
}
{    }{
}
{   }{
}
}
{}{}{}
und

\mavergleichskette
{\vergleichskette
{ B
}
{ = }{ K[Y_1 , \ldots , Y_n ]/ {\mathfrak b}
}
{  }{
}
{    }{
}
{   }{
}
}
{}{}{}
\definitionsverweis {endlich erzeugte}{}{}
$K$-\definitionsverweis {Algebren}{}{.}
Es sei

\mavergleichskettedisp
{\vergleichskette
{  \ell
}
{ =} {  \operatorname{max} \left( m ,\,  n \right)
}
{ } {
}
{ } {
}
{ } {
}
}
{}{}{.}
Zeige, dass man das
$K$-\definitionsverweis {Spektrum}{}{}
vom
\definitionsverweis {Produktring}{}{}

\mathl{A \times B}{} als eine abgeschlossene Menge des
\mathl{{ {\mathbb A}_{ K }^{  \ell +1  } }}{} realisieren kann.





}
{} {}

\inputaufgabe
{}
{





Es sei $X$ ein
\definitionsverweis {topologischer  Raum}{}{,}
der nicht leer und nicht \definitionsverweis {zusammenhängend}{}{}
sei. Zeige, dass es dann eine stetige Abbildung
\mathl{f:X \rightarrow \R}{,}
\mathl{f\neq 0,1}{,}
\zusatzklammer {$\R$ sei mit der metrischen Topologie versehen} {} {}
gibt, die
\definitionsverweis {idempotent}{}{}
im Ring der stetigen Funktionen auf $X$ ist.





}
{} {}

\inputaufgabe
{}
{





Es sei $X$ ein
\definitionsverweis {topologischer Raum}{}{}
mit einer
\definitionsverweis {disjunkten Zerlegung}{}{}

\mavergleichskettedisp
{\vergleichskette
{X
}
{ =} { U \uplus V
}
{ } {
}
{ } {
}
{ } {
}
}
{}{}{}
aus
\definitionsverweis {offenen Teilmengen}{}{}

\mavergleichskette
{\vergleichskette
{U,V
}
{ \subseteq }{ X
}
{  }{
}
{    }{
}
{   }{
}
}
{}{}{.}
Zeige, dass die natürliche Abbildung
\maabbeledisp {} {C(X, \R) } { C(U, \R)  \times C(V, \R)
} { f } { ( f \vert_U, f \vert_V)
} {,}
\definitionsverweis {bijektiv}{}{}
ist.





}
{} {}

\inputaufgabegibtloesung
{}
{





Es sei $R$ ein
\definitionsverweis {kommutativer Ring}{}{}
mit
\definitionsverweis {Reduktion}{}{}
$S$. Zeige, dass die Abbildung, die den
\definitionsverweis {idempotenten Elementen}{}{}
aus $R$ ihre Restklasse in $S$ zuordnet,
\definitionsverweis {injektiv}{}{}
ist.





}
{} {}

\inputaufgabegibtloesung
{}
{





Es sei $R$ ein
\definitionsverweis {kommutativer Ring}{}{}
mit einem Element

\mavergleichskette
{\vergleichskette
{ n
}
{ \in }{ R
}
{  }{
}
{    }{
}
{   }{
}
}
{}{}{}
mit

\mavergleichskette
{\vergleichskette
{ n^2
}
{ = }{ 0
}
{  }{
}
{    }{
}
{   }{
}
}
{}{}{}
in $R$ und sei

\mavergleichskettedisp
{\vergleichskette
{ S
}
{ =} { R/(n)
}
{ } {
}
{ } {
}
{ } {
}
}
{}{}{.}
Zeige, dass es zu jedem
\definitionsverweis {idempotenten Element}{}{}
$e$ aus $S$ ein idempotentes Element aus $R$ gibt, dessen Restklasse gleich $e$ ist.





}
{} {}

\inputaufgabe
{}
{





Es sei $R$ ein
\definitionsverweis {noetherscher}{}{}
\definitionsverweis {kommutativer Ring}{}{}
mit
\definitionsverweis {Reduktion}{}{}
$S$. Zeige, dass es eine Folge von kommutativen Ringen

\mathbed {R_i} {}
 {1 \leq i \leq n} {}
 {} {} {} {,}
und surjektiven
\definitionsverweis {Ringhomomorphismen}{}{}
\maabbdisp {\varphi_i} {R_i } { R_{i+1}
} {}
derart gibt, dass die Gesamtabbildung

\mathdisp {R=R_0 \rightarrow R_1 \rightarrow R_2 \rightarrow  \cdots \rightarrow R_{n-1} \rightarrow R_n = S} {  }

die Reduktionsabbildung ist und jedes $\varphi_i$ der Restklassenhomomorphismus
\maabbdisp {} { R } { R/(x_i)
} {}
zu einem Element

\mavergleichskette
{\vergleichskette
{ x_i
}
{ \in }{ R_i
}
{  }{
}
{    }{
}
{   }{
}
}
{}{}{}
mit

\mavergleichskette
{\vergleichskette
{ x_i^2
}
{ = }{ 0
}
{  }{
}
{    }{
}
{   }{
}
}
{}{}{}
in $R_i$ ist.





}
{} {}

\inputaufgabe
{}
{





Es sei $R$ ein
\definitionsverweis {kommutativer Ring}{}{}
mit
\definitionsverweis {Reduktion}{}{}
$S$. Zeige, dass die Abbildung, die den
\definitionsverweis {idempotenten Elementen}{}{}
aus $R$ ihre Restklasse in $S$ zuordnet,
\definitionsverweis {surjektiv}{}{}
ist.





}
{} {}

Die folgende Aussage ist eine Version des Chinesischen Restsatzes.

\inputaufgabegibtloesung
{}
{





Es sei $R$ ein
\definitionsverweis {kommutativer Ring}{}{}
und es seien

\mathbed {{\mathfrak a}_j} {}
 {j=1  , \ldots , n} {}
 {} {} {} {,}
\definitionsverweis {Ideale}{}{}
mit

\mavergleichskette
{\vergleichskette
{ {\mathfrak a}_i + {\mathfrak a}_j
}
{ = }{ R
}
{  }{
}
{    }{
}
{   }{
}
}
{}{}{}
für alle

\mavergleichskette
{\vergleichskette
{i
}
{ \neq }{j
}
{  }{
}
{    }{
}
{   }{
}
}
{}{}{.}
Zeige, dass

\mavergleichskettedisp
{\vergleichskette
{R/  {\mathfrak a}_1 \cdots  {\mathfrak a}_n
}
{ \cong} { R/{\mathfrak a}_1 \times \cdots \times  R/{\mathfrak a}_n
}
{ } {
}
{ } {
}
{ } {
}
}
{}{}{}
gilt.





}
{} {}





  \zwischenueberschrift{Aufgaben zum Abgeben}

\inputaufgabe
{4}
{





Es sei $R$ ein
\definitionsverweis {Hauptidealbereich}{}{}
mit
\definitionsverweis {Quotientenkörper}{}{}

\mavergleichskette
{\vergleichskette
{Q
}
{ = }{ Q(R)
}
{  }{
}
{    }{
}
{   }{
}
}
{}{}{.}
Zeige, dass jeder Zwischenring

\mathbed {S} {}
 {R \subseteq S \subseteq Q} {}
 {} {} {} {,}
eine
\definitionsverweis {Nenneraufnahme}{}{}
ist.





}
{} {}

\inputaufgabe
{5 (1+2+1+1)}
{





Betrachte die durch

\mavergleichskette
{\vergleichskette
{ Y^2
}
{ = }{ X^3+X^2
}
{  }{
}
{    }{
}
{   }{
}
}
{}{}{}
gegebene Kurve $C$ \zusatzklammer {siehe

Beispiel 6.3
} {} {} und die offene Menge

\mavergleichskette
{\vergleichskette
{ U
}
{ = }{ D(X)
}
{ \subseteq }{ C
}
{    }{
}
{   }{
}
}
{}{}{.}
\aufzaehlungvierabc{Finde eine abgeschlossene Realisierung von $U$ in ${ {\mathbb A}_{ K }^{  3   } }$.
}{Zeige, dass es auch in ${\mathbb A}^{2}_{K}$  eine abgeschlossene Realisierung gibt.
}{Ist $U$ isomorph zu einer offenen Menge der affinen Geraden?
}{Skizziere die Bildkurve unter der Abbildung
\maabbeledisp {} { U } { {\mathbb A}^{2}_{ \R }
} { (x,y) } { \left(  { \frac{   1  }{  x } } , y                     \right)
} {.}
}





}
{} {}

\inputaufgabe
{4}
{





Betrachte zwei parallele Geraden $V$ und das Achsenkreuz $W$. Beschreibe eine möglichst natürliche surjektive Abbildung zwischen $V$ und $W$ (in welche Richtung?), und zwar sowohl geometrisch als auch algebraisch. Gibt es auch eine surjektive polynomiale Abbildung in die andere Richtung?





}
{} {}

\inputaufgabe
{3}
{





Bestimme die
\definitionsverweis {nilpotenten}{}{}
und die
\definitionsverweis {idempotenten}{}{}
Elemente in
\mathl{\Z/( 175)}{.}





}
{} {}

\inputaufgabe
{4}
{





Es sei $K$ ein
\definitionsverweis {algebraisch abgeschlossener Körper}{}{,}
wir betrachten den Durchschnitt der beiden
\definitionsverweis {algebraischen Kurven}{}{}

\mathdisp {V \left(  X^2+Y^2-1  \right)  \text{ und } V \left(  Y-X^2  \right)} { . }

Identifiziere den
\definitionsverweis {Restklassenring}{}{}

\mavergleichskettedisp
{\vergleichskette
{ R
}
{ =} { K[X,Y]/ \left(  X^2+Y^2-1,Y-X^2  \right)
}
{ } {
}
{ } {
}
{ } {
}
}
{}{}{}
mit einem
\definitionsverweis {Produktring}{}{}
und beschreibe die Restklassenabbildung
\mathl{K[X,Y] \rightarrow R}{} mittels dieser Identifizierung. Bestimme Urbilder in
\mathl{K[X,Y]}{} für sämtliche
\definitionsverweis {idempotenten Elemente}{}{}
des Produktringes.





}
{} {}

\inputaufgabe
{6}
{





Es sei $K$ ein Körper und $A$ eine endlichdimensionale,
\definitionsverweis {reduzierte}{}{}
$K$-Algebra. Zeige, dass dann $A$ ein endliches direktes Produkt von endlichen Körpererweiterungen von $K$ ist.





}
{} {Hinweis: Man darf ohne Beweis benutzen, dass es in $A$ nur endlich viele Primideale gibt.}
